\documentclass[./main.tex]{subfiles}
\begin{document}

% Slide # 10
\begin{frame}[t, label=slide10]
        % Title
        \frametitle{Predictive skill of an EWS}

        % Body
        \footnotesize

        \begin{columns}[T]
                
                \column{0.5\textwidth}

                \vspace{-2mm}

                How do I tackle higher dimensions?

                \vspace{5mm}

                \uncover<2->{

                        Consider a saddle-node normal form in $2-$dimensions

                        \begin{equation*}
                                \begin{cases}
                                        dx_{t} = -\mu(t) - x_{t}^{2} + \sigma dW_{t}\,, \\
                                        dy_{t} = -y_{t} + \sigma dW_{t}\,.
                                \end{cases}
                        \end{equation*}
                }

                \vspace{4mm}

                \uncover<3->{
                        Choose a scalar observable $\psi(x_{t},y_{t}) =: \psi_{t}$ that exhibits tipping.
                }

                \vspace{5mm}

                \uncover<4->{
                        The predictive skill of the increase in variance EWS depends on how well the designed observable capture CSD.
                }

                \hfill

                \column{0.5\textwidth}

                \vspace{-7mm}

                \begin{figure}[H]
                        \centering 
                        \includegraphics<3->[keepaspectratio, width=0.85\textwidth]{./figures/slide_10_1.png}%
                \end{figure}

                \vspace{-6mm}

                \begin{figure}[H]
                        \centering 
                        \includegraphics<4->[keepaspectratio, width=0.85\textwidth]{./figures/slide_10_2.png}%
                \end{figure}

       \end{columns}

\end{frame}

\end{document}
